\chapter{元件層級設計}
\label{ch:device}
\section{低電壓元件最佳化}
針對標稱高效能操作最佳化的元件，不一定能在 near-threshold 區達到最低能量。目標汲極電壓降低後，短通道控制、接面電容、串聯電阻與 subthreshold swing 的平衡也會改變。Dreslinski 等人回顧了調整通道摻雜及元件幾何的研究 \cite[Sec.~V-A]{dreslinski2010}；其改善幅度屬於特定製程結果，並非任意製程皆應降低摻雜的通則。

較陡的 subthreshold 特性，可在給定閘極電壓範圍內增加導通與截止狀態的區別。較低的接面或佈線電容可減少電荷搬移，較低的串聯電阻則可能提高可用驅動能力。這些效應必須共同評估：提高電流的幾何變更可能增加電容，而有利於低電壓的調整，也可能犧牲較高電壓下的峰值效能。

重要的評估標準，是代表性電路的能量延遲及穩健性，而非孤立的元件指標。截止電流較低，但有效驅動電流也大幅下降的電晶體，可能因執行時間延長而增加漏電能量。同樣地，標稱 $\Ion/\Ioff$ 比值改善，也不能單獨保證 SRAM 寫入能力，因為寫入取決於元件相對強度及實際偏壓。

實體實作也須納入同一分析。互連、時脈樹與電源佈線，可能消耗閘級估計所預測的時序餘裕。較低的目標電壓，可能需要更強的電源分配或額外去耦來控制電壓下陷。這些措施帶來面積與電容成本，其價值在於能否支援額外且可靠的降壓。

\section{電晶體尺寸設計}
增加電晶體尺寸可提高驅動能力或降低部分失配效應，但會增加電容、面積，且常伴隨漏電上升。最小尺寸元件雖有節能吸引力，卻可能不滿足功能與時序要求。因此尺寸選擇是受限制的最佳化問題，而不是所有低電壓元件一律放大或縮小。

Calhoun 等人在 0.18-$\mu$m FIR 測試晶片中，比較傳統實作與降低最低操作電壓的修改版本 \cite{calhoun2005}。兩個濾波器都可在很低的電壓操作，但修改版本在典型條件下約有 50\% 的額外能量成本。原文將其中相當部分歸因於受限的單元集合及其合成結果；在所比較的 worst-case corner 下，能量改善則低於 10\%。

此結果對低電壓元件庫的選擇尤其重要。移除複雜或較弱的單元，可改善功能可行性，卻可能迫使合成使用更多邏輯閘或更長路徑。只放大少數關鍵單元，可能比全面調整尺寸更有效率，但仍須檢查變異。唯有完整設計滿足能量與吞吐量目標，較低的功能電壓才有價值。

最低能量附近的曲線可能相當平坦。若稍高的供應電壓能允許更小、更簡單且良率足夠的實作，其表現可能優於大幅調整尺寸後、在較低電壓操作的電路。因此應比較完整實作的可行能量曲線，而不只比較最低成功測試電壓。

\section{臨界電壓配置}
降低臨界電壓可在降壓後恢復速度，但通常增加漏電；高臨界電壓單元則可降低具有足夠時序餘裕之路徑的漏電。多臨界電壓設計依照各路徑角色配置單元，再於目標供應電壓重新評估時序與能量 \cite{chandrakasan1992,gonzalez1997}。由於路徑相對延遲會變化，標稱電壓下適合的配置，在 near-threshold 區可能需要調整。

多臨界電壓單元配置，應與常以 MTCMOS 指稱的 power gating 結構區分。高臨界電壓睡眠元件會引入虛擬電源、主動模式壓降、喚醒行為與狀態保持需求。當核心本來就缺乏電壓裕量時，這些成本尤其重要。本報告僅在 power gating 改變漏電與可行電壓策略的範圍內討論它，並不將其作為主題。

降低臨界電壓的一階效益容易被高估。在弱反轉區，改變臨界電壓可能同時增加有效電流與漏電流；較快執行與較大漏電，可能在任務能量上部分互相抵消。因此合理比較應評估 $\VDD I_{\mathrm{leak}}T_{\mathrm{task}}$，而非假定恢復速度一定能節能。

\section{基體偏壓控制}
當製程允許有效且安全的基體控制時， body bias可改變臨界電壓，進而調整速度與漏電。傳統體矽元件可用下式說明其相依性：
\begin{equation}
\VTH(V_{\mathrm{SB}})=V_{\mathrm{th0}}+
\gamma\left(\sqrt{2\phi_F+V_{\mathrm{SB}}}-\sqrt{2\phi_F}\right).
\end{equation}
此簡化 NMOS 關係在其假設範圍內，說明 body effect 的變化方向；它不是所有全空乏或多閘極製程都適用的通用模型。

獨立控制 p 型與 n 型元件的基體，可協助平衡驅動能力，對依賴元件比例的低電壓功能很有價值。 near-threshold 綜述討論以 adaptive body bias 因應全域及局部變異 \cite[Sec.~V-B]{dreslinski2010}。然而，全域偏壓無法消除所有隨機局部失配，偏壓網路本身也消耗面積與能量。

 forward body bias 可用漏電增加換取速度， reverse body bias 則可能降低待機漏電但犧牲效能。可用範圍、井隔離、接面限制與敏感度都取決於製程。對 voltage scaling 設計，應判斷 body bias 的總成本是否小於更佳供應電壓設定帶來的節能，而非僅因電路操作於 near-threshold 區就加入此技術。

\clearpage
\section{元件設計策略}
\begin{table}[H]
\centering\tablebegin
\caption{支援供應電壓降低的元件與實體設計措施。}
\begin{tabularx}{\textwidth}{@{}p{0.25\textwidth}YY@{}}
\toprule 措施 & 預期效益 & 成本或待驗證條件 \\\midrule
低電壓元件最佳化 & 改善有效驅動、擺幅或電容。 & 製程特定的漏電、可靠度與高電壓行為。 \\
選擇性尺寸調整 & 恢復關鍵路徑時序或穩健性。 & 額外電容與漏電。 \\
臨界電壓配置 & 在需要之處提供速度。 & 關鍵路徑變更與漏電能量。 \\
 adaptive body bias & 追蹤變異或平衡元件驅動。 & 製程支援、偏壓限制與控制成本。 \\
改善電源網路 & 降低局部壓降、恢復電壓裕量。 & 佈線、去耦與面積成本。 \\\bottomrule
\end{tabularx}
\end{table}

這些技術分別影響最佳化問題中的不同項目。唯有其組合擴大可靠且節能的操作區域，才有採用理由。若改善單一指標，卻使整個系統偏離該區域，便不能視為成功的 voltage scaling 結果。

